\chapter{Кратности общего неприводимого следового носителя}
\label{ch:v7-common-irreducible-trace-multiplicity-gate}

\section{Развилка одного следа}

Текущие кривизны естественно живут на двух носителях размерностей
\begin{equation}
 \dim\mathcal K_E=22,
 \qquad
 \dim\mathcal K_V=11+10=21.
 \label{eq:v7-common-trace-carrier-dimensions}
\end{equation}
Если сохранять их блочно-диагонально, алгебра равна
\begin{equation}
 M_{22}(\mathbb C)\oplus M_{21}(\mathbb C),
 \qquad
 \tau=w\tau_{22}+(1-w)\tau_{21}.
 \label{eq:v7-common-trace-central-fork}
\end{equation}
После общей нормировки остаётся свободный вес $w$. Переход к простому
$M_{43}(\mathbb C)$ удаляет центр, но одновременно вводит полный
внедиагональный угол
\begin{equation}
 \operatorname{Hom}(\mathbb C^{21},\mathbb C^{22}),
 \qquad \dim_{\mathbb C}=462,
 \label{eq:v7-common-trace-offdiagonal-cost}
\end{equation}
то есть новый класс коннекторов и новое действие.

Для одной простой матричной алгебры нормированный след единственен. Однако
равномерная кратность представления не создаёт коэффициент:
\begin{equation}
 \frac1{nk}\Tr_{nk}(X\otimes I_k)=\frac1n\Tr_nX.
 \label{eq:v7-common-trace-multiplicity-cancellation}
\end{equation}
В текущей конструкции оба блока представлены по одному разу, поэтому общий
полный след снова даёт $\beta=1$.

\section{Почему резонанс \texorpdfstring{$22\to21$}{22 -> 21} ещё не является оператором}

Можно было бы попытаться считать два блока квадратами одного прямоугольного
оператора. Но для любого $B$ выполняется
\begin{equation}
 \rank(B^*B)=\rank(BB^*).
 \label{eq:v7-common-trace-rank-pairing}
\end{equation}
Текущий рёберный опорный квадрат имеет ранг $22$, тогда как два физических
Gram-конца единичного корневого фона имеют общий ранг $20$. После дополнения
нулём евклидова невязка их спектров равна
\begin{equation}
 \|\lambda_E-\lambda_V\|_2=4.
 \label{eq:v7-common-trace-spectral-mismatch}
\end{equation}
Следовательно, размеры $22$ и $21$ образуют полезную индексную подсказку, но
текущие кривизны не являются двумя концами одного Hodge--Dirac квадрата.

\section{Привлекательное отношение \texorpdfstring{$11/21$}{11/21}}

Нормированный след физического source-проектора равен
\begin{equation}
 \tau_{21}(P_{11})=\frac{11}{21}<\frac8{15}.
 \label{eq:v7-common-trace-source-fraction}
\end{equation}
Если просто умножить всю физическую норму на это число, гессиан действительно
проходит:
\begin{equation}
 \beta=\frac{11}{21}:qquad
 (n_-,n_0,n_+)=(7,0,20),
 \qquad \lambda_{\rm heavy}^{\min}=\frac4{35}.
 \label{eq:v7-common-trace-source-fraction-pass}
\end{equation}
Это наиболее близкая числовая зацепка текущей базы.

Но настоящий оператор $P_{11}$ внутри следа не умножает всю двустороннюю
Gram-кривизну. Он оставляет только source-конец. Аналогично $P_{10}$ оставляет
только target-конец. Явный аудит даёт
\begin{equation}
 (n_-,n_0,n_+)_{P_{11}}=(8,0,19),
 \qquad
 (n_-,n_0,n_+)_{P_{10}}=(17,0,10).
 \label{eq:v7-common-trace-corner-insertion-signatures}
\end{equation}
Оба честных проекторных следа проваливаются. Операция
$\tau_{21}(P_{11})\mathcal S_V$ является произведением уже вычисленных
скаляров, а не одним следом кривизны.

\begin{theorem}[No-go простого следа и условный размерностный резонанс]
Ни общий полный след одной копии, ни нормированные corner-следы, ни текущая
кратность представлений не выводят коэффициент в окне
$0\le\beta<8/15$. Число $11/21$ даёт правильный гессиан только как
постфактумный множитель всего физического сектора. Его честная вставка как
source-проектора выбирает один Gram-конец и не проходит.
\end{theorem}

\begin{nogocriterion}[Запрет умножения на след проектора]
Нельзя заменять $\Tr(P\mathfrak m_V^2)$ произведением
$\tau(P)\Tr(\mathfrak m_V^2)$. Эти выражения совпадают лишь при специальной
изотропии кривизны, которой нет на полном пространстве вариаций. Для спасения
$11/21$ требуется заранее заданный передаточный оператор или условное
ожидание, действующее на саму кривизну.
\end{nogocriterion}

\begin{protocol}
\begin{enumerate}
 \item блочно-диагональный центр: \textbf{два};
 \item свободных нормированных весов: \textbf{один};
 \item простой completion: \textbf{$M_{43}(\mathbb C)$};
 \item новых комплексных коннекторов: \textbf{$462$};
 \item равномерная кратность меняет нормированный след: \textbf{нет};
 \item общий однокопийный вес: \textbf{$\beta=1$};
 \item ранги предполагаемой Hodge-пары: \textbf{$22$ и $20$};
 \item числовой кандидат $11/21$: \textbf{проходит условно};
 \item честный source-corner: \textbf{$(8,0,19)$, провал};
 \item честный target-corner: \textbf{$(17,0,10)$, провал};
 \item итог: \textbf{простой след и текущие кратности закрыты};
 \item следующий гейт: \textbf{инцидентное условное ожидание или марковский вес}.
\end{enumerate}
\end{protocol}

Машинный сертификат сохранён в
\path{s2t_v7_common_irreducible_trace_multiplicity_gate_results.json}.

\begin{thebibliography}{9}
\bibitem{v7-common-trace-krajewski}
T.~Krajewski,
\emph{Classification of Finite Spectral Triples},
J. Geom. Phys. 28 (1998), 1--30; arXiv:hep-th/9701081.
\bibitem{v7-common-trace-af-lift}
T.~Masson, G.~Nieuviarts,
\emph{Lifting Bratteli Diagrams between Krajewski Diagrams: Spectral
Triples, Spectral Actions, and AF Algebras}, arXiv:2207.04466.
\end{thebibliography}